\chapter{Reflexões sobre a Benevolência}

Que eu permaneça em bem-estar\\
Livre de aflições,\\
Livre da hostilidade,\\
Livre da má vontade,\\
Livre da ansiedade,\\
E que eu possa manter o bem-estar em mim mesmo.\\
% REVISÃO: concordância — referindo-se a "todos" (plural), deveria ser "Livres da hostilidade / Livres da má vontade / Livres da ansiedade" (no plural).
% REVISÃO (ortografia): "ações"/"atos" estão na grafia do Acordo de 1990; o grosso do livro usa a grafia pré-Acordo ("acções"/"actos"). Uniformizar.
Que todos permaneçam em bem-estar,\\
Livre da hostilidade,\\
Livre da má vontade,\\
Livre da ansiedade,\\
E que possam manter o bem-estar em si mesmos.\\
Que todos os seres sejam libertados de todo o sofrimento\\
E que não se separem da boa sorte que alcançaram.\\
Quando agem com intenção\\
Todos os seres são donos das suas ações,\\
E herdam os seus resultados.\\
O seu futuro nasce dessas ações,\\
E os seus resultados serão o seu lar.\\
Todas as ações com intenção,\\
Sejam elas hábeis ou prejudiciais;\\
De tais atos, eles serão os herdeiros.
